\section{Datos del entorno y de la instalación}

La práctica se realizó en una estación con Debian GNU/Linux 13 (trixie), de
arquitectura x86\_64. El núcleo registrado fue
\texttt{6.12.107+deb13-amd64}. La instalación de DBeaver correspondió a la
edición Community 26.2.0.

\begin{table}[H]
  \centering
  \rowcolors{2}{unamblue!6}{white}
  \begin{tabularx}{\textwidth}{>{\bfseries}lX}
    \toprule
    Componente & Versión registrada \\
    \midrule
    Sistema operativo & Debian GNU/Linux 13 (trixie), x86\_64 \\
    Núcleo & \texttt{6.12.107+deb13-amd64} \\
    Docker Engine & 29.8.0 \\
    Docker Compose & 5.5.1 \\
    Java & OpenJDK 21.0.12 \\
    DBeaver Community & 26.2.0 \\
    PostgreSQL & 18.6, dentro de la imagen \texttt{postgres:18} \\
    \bottomrule
  \end{tabularx}
  \caption{Resumen del entorno utilizado durante la práctica.}
\end{table}

\evidencia{01_evidencia.png}{Identificación del sistema operativo, Docker, Java y DBeaver.}

No se registró un tiempo exacto de trabajo. La práctica incluyó la validación
del entorno, la ejecución de Docker, el despliegue del contenedor de
PostgreSQL y la conexión desde DBeaver.
